% ============================================================
% RESOLVE Swiss — statistical analysis plan
% Target: Brazilian Journal of Physical Therapy, "Clinical Trial Protocol"
% (the SAP of the Australian RESOLVE trial appeared there: Bagg et al. 2021)
%
% STATUS: frequentist analysis only. The Bayesian interim analysis is
% deliberately not covered yet; see the placeholder in section
% "Interim analysis" and TODO.md.
%
% Numbers in the Sample size section are produced by R/sample_size.R.
% Build: latexmk -pdf paper.tex
% ============================================================
\documentclass[11pt,a4paper]{article}

\usepackage[utf8]{inputenc}
\usepackage[T1]{fontenc}
\usepackage{lmodern}
\usepackage[english]{babel}
\usepackage{microtype}
\usepackage[a4paper,margin=2.6cm]{geometry}
\usepackage{setspace}
\usepackage{amsmath}
\usepackage{booktabs}
\usepackage{tabularx}
\usepackage{array}
\usepackage{enumitem}
\usepackage{titlesec}
\usepackage{xcolor}
\usepackage[numbers,sort&compress]{natbib}
\usepackage[hidelinks]{hyperref}
\usepackage[mathlines]{lineno}

\titleformat{\section}{\large\bfseries}{\thesection}{0.6em}{}
\titleformat{\subsection}{\normalsize\bfseries}{\thesubsection}{0.6em}{}
\titleformat{\subsubsection}{\normalsize\itshape}{\thesubsubsection}{0.6em}{}
\setlength{\parskip}{0.4em}
\setlength{\parindent}{0pt}
\renewcommand{\arraystretch}{1.2}

% open items that still need input from the trial team
\newcommand{\open}[1]{\textcolor{red}{[\textsc{open}: #1]}}

\doublespacing
\pagestyle{plain}

\title{\bfseries The RESOLVE Swiss trial for people with chronic low back pain:\\
statistical analysis plan for a cluster randomised controlled trial}
\author{Jürgen Degenfellner\thanks{Corresponding author.
Institute of Physiotherapy, ZHAW Zurich University of Applied Sciences,
Katharina-Sulzer-Platz 9, 8400 Winterthur, Switzerland.
E-mail: \texttt{jdegenfellner@gmail.com}}
\ \open{co-author list and affiliations}}
\date{}

\begin{document}
\linenumbers
\maketitle

% ============================================================
\section*{Abstract}

\textbf{Background.} Statistical analysis plans describe the planned data
management and analysis of a clinical trial before the data are seen, and so
protect the eventual results from data-driven analytical choices. This paper
reports the statistical analysis plan for the RESOLVE Swiss trial, which
evaluates graded sensorimotor retraining combined with pain neuroscience
education against usual physiotherapy for people with chronic low back pain.

\textbf{Design and analysis.} RESOLVE Swiss is a two-arm parallel-group
\emph{cluster} randomised controlled trial in Swiss physiotherapy practices;
the practice is the unit of randomisation and the patient is the unit of
analysis. The primary outcome is back-specific disability, measured with the
Roland--Morris Disability Questionnaire at 18 weeks after randomisation. The
primary analysis is a linear mixed model with a random intercept for practice
and adjustment for the baseline value of the outcome, estimating the
participant-average difference between arms with Kenward--Roger degrees of
freedom. Because the trial randomises a small number of practices, and because
mixed models are known to produce anti-conservative inference in that setting,
we prespecify two further analyses of the primary outcome: a cluster-level
(two-stage) analysis and an exact randomisation test over the finite set of
possible allocations. We report the intra-cluster correlation with its
uncertainty, the analyses of secondary outcomes, the handling of missing data,
non-adherence and contamination, and the planned reporting.

\textbf{Conclusion.} Prespecifying the analysis reduces the scope for bias in
the results of RESOLVE Swiss and makes explicit how the small number of
clusters is dealt with.

\textbf{Trial registration.} \open{registry and registration number}

\vspace{0.4em}
\noindent\textbf{Keywords:} Low back pain; Cluster randomised trial;
Statistical analysis plan; Mixed models; Physical therapy.

% ============================================================
\section{Introduction}

\subsection{Background and rationale}

Chronic low back pain is among the largest contributors to years lived with
disability worldwide, and the effects of the treatments physiotherapists
currently offer are, on average, modest. The Australian RESOLVE trial tested
whether a treatment aimed at the central nervous system representation of the
back --- graded sensorimotor precision training preceded by pain neuroscience
education --- reduces pain more than a sham intervention, and reported a small
but consistent benefit at 18 weeks \citep{bagg2022}. An earlier and much
smaller Swiss pilot compared a related programme with usual physiotherapy and
pointed in the same direction without being able to resolve the question
\citep{waelti2015}.

Two things are unresolved. First, RESOLVE compared the intervention with a
sham; the clinically relevant comparison for Swiss practice is against usual
physiotherapy as it is actually delivered. Second, the intervention has so far
been delivered in a trial setting by trial therapists. RESOLVE Swiss therefore
delivers it in ordinary physiotherapy practices, by the therapists who work
there. That choice makes the trial pragmatic, and it forces the design to be a
\emph{cluster} randomised trial: once a practice has been trained in the
intervention, its therapists cannot un-learn it for the next patient, so
allocation has to happen at the level of the practice rather than the patient
\citep{campbell2012}.

Cluster randomisation changes the statistics in ways that are easy to get
wrong, and the difficulty is sharpened here because the number of practices is
small. This analysis plan is written before the analysis of outcome data and
follows the recommendations for the content of statistical analysis plans in
clinical trials \citep{gamble2017}, together with the CONSORT extension for
cluster randomised trials \citep{campbell2012}. It reports the planned
frequentist analysis of the primary and secondary outcomes.
\open{decide whether the Bayesian interim analysis is reported in this paper or
in a companion paper; see section~\ref{sec:interim}}

% ============================================================
\section{Methods}

\subsection{Trial design}

RESOLVE Swiss is a two-arm, parallel-group, cluster randomised controlled
trial with 1:1 intended allocation at the level of the physiotherapy practice.
Practices are the clusters; patients treated within a practice are the units of
observation and of analysis. Outcome data are collected at baseline and at
18 weeks after randomisation \open{confirm the full follow-up schedule; the
Australian trial also collected 26 and 52 weeks}.

At the time of writing, \open{confirm} 5 practices are allocated to the
intervention arm and 9 to the control arm. The consequences of that split, and
of any change to it, are quantified in section~\ref{sec:samplesize}.

Ethical approval was granted by \open{ethics committee and approval number}.

\subsection{Eligibility}

A cluster randomised trial has two eligibility layers, and both are reported
here because selection can operate at either \citep{campbell2012}.

\emph{Practices.} \open{practice-level inclusion criteria: e.g. minimum patient
volume, number of therapists, willingness to be randomised to either arm,
completion of the intervention training}

\emph{Participants.} \open{participant-level criteria; from the sample size
work the trial requires a baseline RMDQ of at least 7 points, please confirm
and complete the remaining criteria}

\subsection{Interventions}

\open{brief description of the intervention and of usual physiotherapy as the
comparator, with a reference to the protocol; note the number of sessions and
the training the intervention therapists received}

Usual physiotherapy is expected to be heterogeneous between practices. What
control practices actually deliver will be recorded and described, because in a
pragmatic trial the comparator is part of the finding rather than a nuisance.

\subsection{Randomisation and allocation}

\open{describe how the practices were allocated, who generated the sequence,
and whether allocation was simple, stratified, or covariate-constrained}

Two points matter for the analysis. First, any variable used to stratify or
constrain the allocation must be adjusted for in the analysis, otherwise the
standard errors are wrong \citep{kahan2012}. Second, if practices were recruited
in waves and later waves were allocated non-randomly, wave and treatment arm
would be confounded; recruitment wave will therefore be recorded for every
practice and, if practices entered in distinct waves, carried in the model as a
covariate.

Participants cannot be blinded to the treatment they receive. Outcome
questionnaires are self-reported. \open{state who is blinded --- data manager,
statistician --- and confirm that the statistician will be blinded to arm
labels until the primary analysis is complete}

\subsection{Recruitment of participants after cluster allocation}

Because practices are randomised before their patients are recruited, and
because the recruiting therapists know which arm their practice is in, patient
selection can differ between arms. This identification or recruitment bias is
the characteristic threat to cluster randomised trials and is not present in
individually randomised trials \citep{campbell2012, eldridge2009}. We will
therefore report, by arm and by practice, the number of patients screened,
found eligible, approached, and consenting, and we will compare baseline
characteristics between arms with this threat explicitly in mind. Baseline
imbalance in a cluster randomised trial with few clusters is expected by chance
and is not by itself evidence of bias; it is, however, the main signal that
recruitment differed between arms, and it will be interpreted alongside the
screening counts rather than by significance testing
\citep{campbell2012, senn1994}.

\subsection{Sample size}
\label{sec:samplesize}

The target difference is 2 points on the Roland--Morris Disability
Questionnaire (RMDQ, range 0--24), the smallest between-group difference the
investigators consider clinically worthwhile in this population. The between
participant standard deviation of the RMDQ at 18 weeks is taken from the
Australian trial (intervention 3.6, SD 4.6; control 6.4, SD 5.8
\citep{bagg2022}), giving a pooled SD of 5.2 points. With a two-sided
$\alpha = 0.05$ and 80\% power, an individually randomised trial would need
108 participants per arm; adjusting the 18-week score for its baseline value,
with an assumed baseline-to-follow-up correlation of 0.6, reduces this to 69
per arm.

Clustering inflates this requirement by the design effect. Allowing for unequal
practice sizes \citep{eldridge2006},
\begin{equation}
\label{eq:de}
\mathrm{DE} = 1 + \left[ \left( \mathrm{cv}^2\,\frac{k-1}{k} + 1 \right) \bar m - 1 \right]\rho ,
\end{equation}
where $\bar m$ is the mean number of analysed participants per practice, $k$ the
number of practices, $\rho$ the intra-cluster correlation and $\mathrm{cv}$ the
coefficient of variation of practice size. Intra-cluster correlations reported
for primary care outcomes are usually small, with a median near 0.01
\citep{adams2004}. Unequal cluster size can be ignored when $\mathrm{cv} < 0.23$,
but for trials randomising general practices $\mathrm{cv}$ is typically around
0.65 \citep{eldridge2006}; physiotherapy practices are the closest analogue we
have, so we plan on $\mathrm{cv} = 0.65$ and report the more optimistic 0.4 as
well. With $\bar m = 14$, $k = 14$ and $\mathrm{cv} = 0.65$, the design effect
is 1.19 at $\rho = 0.01$, 1.56 at $\rho = 0.03$ and 1.93 at $\rho = 0.05$, so
the requirement is 82, 107 and 133 participants per arm respectively. The trial
therefore aims for 100 analysed participants per arm, which corresponds to 110
recruited per arm allowing 10\% loss to follow-up (Table~\ref{tab:ss}); this
covers an intra-cluster correlation up to about 0.03.

Two approximations in this calculation should be named. The variance reduction
from baseline adjustment is applied as the individual-level factor $1-r^2$ and
the clustering inflation as \eqref{eq:de}; a formula that handles baseline
adjustment inside a cluster randomised design, using the individual-level and
the cluster-level baseline correlations separately, is available and gives a
slightly smaller requirement when the cluster-level correlation is high
\citep{teerenstra2012}. Our version is therefore mildly conservative. Second,
the intra-cluster correlation is assumed, not known; the calculation is
reported across a range of values rather than at a single one, which is what
sensitivity to this assumption requires \citep{hemming2023}.

Three features of the design as it currently stands should be stated plainly
rather than buried. First, with 14 practices the between-cluster component of
the variance rests on 14 numbers, and the relevant reference distribution has
roughly 12 degrees of freedom; power is correspondingly lower than the design
effect alone suggests. Holding the analysed sample at 200 participants and
using a $t$ reference distribution on $k_1 + k_2 - 2$ degrees of freedom, the
power to detect a 2-point difference is 77\% at $\rho = 0.01$, 65\% at
$\rho = 0.03$ and 56\% at $\rho = 0.05$.

Second, rebalancing the existing practices from 5 versus 9 to 7 versus 7
recovers only about 3 percentage points of power at every intra-cluster
correlation considered, because the total analysed sample is what it is. Equal
allocation of clusters is the efficient choice when the intra-cluster
correlation and the total variance are the same in both arms, and a deliberate
imbalance is justified only when they are not \citep{copas2021}; we have no
reason to expect them to differ here, so 7 versus 7 is preferable, but the
efficiency argument for changing the split is weak.

Third --- and this is the binding constraint --- \emph{the number of practices
is too small}. Five clusters in one arm sits at the bottom of the range in
which the recommended analyses have been evaluated \citep{leyrat2018}, and a
minimum of about ten clusters per arm has been recommended for cluster
randomised trials generally, with the practical rule of adding two to three
clusters per arm to any sample size derived from a normal approximation
\citep{vanbreukelen2018}. Few clusters also raise the probability of chance
imbalance on practice characteristics and reduce the number of analyses
available \citep{taljaard2016}. Adding practices while holding practice size at
14 participants raises power from 62\% with 5 per arm to 80\% with 7 per arm
and 90\% with 9 per arm at $\rho = 0.01$ (Table~\ref{tab:power}) --- a much
better return than adding patients to the practices already enrolled, because
increasing cluster size at a fixed number of clusters has strongly diminishing
returns \citep{hemming2017}. If additional practices are recruited, they must
be randomised as a batch rather than assigned to the intervention arm, so that
the new wave contains its own randomised contrast
\citep{ivers2012, moulton2004}.

All sample size and power computations are reproducible from the script
\texttt{R/sample\_size.R} that accompanies this paper.

\subsection{Data integrity and missing data}

\open{describe the data capture system (REDCap), double entry or electronic
capture, and the data-cleaning steps}

Data will be inspected for implausible values, range violations and duplicate
records before the analysis begins, and the cleaned data set will be locked
before any comparison between arms is made. We will report the amount and the
pattern of missing outcome data by arm and by practice, since in a cluster
trial an entire practice can be lost, which is a different problem from
individual dropout.

\subsection{Analytic principles}

\begin{itemize}[leftmargin=1.4em, itemsep=2pt]
\item Participants will be analysed in the arm to which their practice was
allocated, regardless of the treatment actually delivered (intention to treat).
\item All treatment effects will be reported as point estimates with 95\%
confidence intervals. Two-sided tests at a nominal $\alpha$ of 0.05 accompany
the primary and secondary analyses, but the interpretation will rest on the
estimate and its interval rather than on whether a threshold was crossed
\citep{wasserstein2019, amrhein2019}.
\item Confidence intervals and $p$ values will not be adjusted for
multiplicity. There is one primary outcome at one time point; secondary
outcomes are explicitly labelled as secondary and will be interpreted as such
\citep{gamble2017}.
\item Every analysis accounts for clustering by practice. No analysis will
treat patients from the same practice as independent.
\item Analyses will be carried out in R \citep{rcore}, with mixed models fitted
using \texttt{lme4} \citep{bates2015} and small-sample inference obtained with
\texttt{lmerTest} and \texttt{pbkrtest} \citep{kuznetsova2017, halekoh2014}.
Analysis code will be published with the results.
\end{itemize}

\subsection{Outcome definitions}

\subsubsection{Primary outcome}

Back-specific disability at 18 weeks after randomisation, measured with the
Roland--Morris Disability Questionnaire, a 24-item instrument scored from 0
(no disability) to 24 \citep{roland1983, chiarotto2016}.

\subsubsection{Secondary outcomes}

\open{list secondary outcomes and their instruments and time points, e.g. pain
intensity (NRS), quality of life (EQ-5D-5L), global perceived effect, adverse
events, credibility/expectancy, and any process measures}

% ============================================================
\section{Analysis}

\subsection{Baseline description}

We will describe the sample at two levels, as the CONSORT cluster extension
requires \citep{campbell2012}. At the practice level (Table~\ref{tab:base}, panel
A) we will report the number of practices per arm, the number of therapists,
practice size, urban or rural setting, and the number of participants
contributed. At the participant level (panel B) we will report means and
standard deviations for approximately normally distributed continuous
variables, medians and interquartile ranges otherwise, and counts with
percentages for categorical variables. Baseline characteristics will not be
tested for statistical significance between arms: the null hypothesis of no
difference is true by construction under randomisation, so such a test answers
no question anyone is asking \citep{altman1985, senn1994}.

\subsection{Primary outcome}

\subsubsection{Estimand}

The quantity of interest is the \emph{participant-average} difference in mean
RMDQ at 18 weeks between the two arms: the difference that would be seen if
every patient of this kind in these practices were treated one way rather than
the other, with each patient counting equally. In a cluster randomised trial
with unequally sized clusters this is not the same quantity as the
cluster-average effect, in which each practice counts equally, and different
analyses target different quantities \citep{kahan2023}. The choice is made here,
at the design stage, as the consensus extension of the ICH E9(R1) estimand
addendum to cluster randomised trials requires \citep{kahan2026}. We prespecify
the participant-average effect and report the cluster-average effect alongside
it.

One consequence deserves stating, because it constrains how the primary model
below may be read. A mixed model with a random intercept weights clusters by
their inverse variance, which is neither participant-average nor cluster-average
weighting; when cluster size is informative --- that is, when outcomes or
treatment effects vary systematically with practice size --- such a model
targets neither estimand cleanly \citep{kahan2023}. The cluster-level analyses
specified in section~\ref{sec:smallk} recover each estimand explicitly through
their weighting, which is a second reason to report them rather than treating
them as a mere robustness check. We will report the distribution of practice
sizes so that readers can judge how much this matters here.

\subsubsection{Primary model}

Let $y_{ij}$ denote the 18-week RMDQ score of participant $i$ in practice $j$,
$y_{ij}^{0}$ the same participant's baseline score, and $x_j = 1$ if practice
$j$ was allocated to the intervention and $0$ otherwise. The primary analysis
is the linear mixed model
\begin{equation}
\label{eq:primary}
y_{ij} = \beta_0 + \beta_1 x_j + \beta_2\, y_{ij}^{0} + u_j + \varepsilon_{ij},
\qquad u_j \sim N(0,\tau^2), \quad \varepsilon_{ij} \sim N(0,\sigma^2),
\end{equation}
where $u_j$ is a random intercept for practice and $\beta_1$ is the treatment
effect. Analysis of covariance on the follow-up score is preferred to analysing
change scores: it is unbiased under randomisation, it is more precise whenever
baseline and follow-up are positively correlated, and it does not inherit the
regression-to-the-mean behaviour of change scores \citep{vickers2001,
twisk2018}; for cluster randomised trials specifically, analysis of covariance
and the equivalent constrained-baseline longitudinal formulation are the
recommended ways to use a pre-randomisation measurement of the outcome
\citep{hooper2018}. The mixed-model analysis of covariance estimator is
moreover consistent under misspecification of the working model, and with few
clusters the main-effects form used here is preferable to one carrying
treatment-by-covariate interactions, which spends degrees of freedom and
requires a robust variance estimator \citep{wang2026}.

Covariate adjustment in a cluster randomised trial is not automatically
beneficial, unlike in an individually randomised one: each cluster-level
covariate costs a degree of freedom, so that the degrees of freedom available
for the treatment effect are approximately the number of practices minus two
minus the number of cluster-level covariates \citep{hemming2025}. With 14
practices the budget is 12, and we therefore adjust for the individual baseline
score --- an individual-level covariate, which does not consume cluster-level
degrees of freedom --- and for nothing else except the variables used to
stratify or constrain the randomisation, which must be adjusted for or the type
I error rate is not maintained \citep{kahan2012, hemming2023}.

The model will be fitted by restricted maximum likelihood. Confidence intervals
and $p$ values for $\beta_1$ will use the Kenward--Roger approximation to the
degrees of freedom and the associated small-sample adjustment to the covariance
matrix of the fixed effects \citep{kenward1997}, which is the standard remedy
for the anti-conservatism of naive mixed-model inference when the number of
clusters is small \citep{leyrat2018, kahan2016}.

\subsubsection{Inference with a small number of clusters}
\label{sec:smallk}

Uncorrected mixed models and generalised estimating equations have inflated
type I error rates when the number of clusters is small, and the inflation is
substantial: in a review and reanalysis of published cluster randomised trials,
a large share of trials with few or moderately many clusters had used methods
that do not maintain the nominal error rate \citep{kahan2016, taljaard2016}. In
a simulation across designs with 4 to 40 clusters, Leyrat and colleagues found
that for continuous outcomes the Satterthwaite and Kenward--Roger corrections
to a mixed model, and variance-weighted cluster-level analysis, did maintain
the nominal rate, whereas uncorrected approaches did not \citep{leyrat2018}.
Model \eqref{eq:primary} with Kenward--Roger degrees of freedom is therefore a
defensible primary analysis for the RMDQ, and it is what the current guidance
for cluster trials with fewer than about 40 clusters recommends
\citep{hemming2023, billot2024}. It is nonetheless an asymptotic argument
resting on 14 numbers, so we prespecify two further analyses of the primary
outcome, reported alongside the model-based result rather than instead of it.
Cluster-level analyses, by contrast, need no such correction at all, which is
why they are the natural companion when the number of clusters is very small
\citep{billot2024}.

\emph{Cluster-level analysis.} We will regress $y_{ij}$ on $y_{ij}^{0}$ within
practices, take the practice-level mean residual as a single summary per
practice, and compare the 14 summaries between arms with a two-sample $t$ test
on $k_1 + k_2 - 2 = 12$ degrees of freedom, weighting practices by their
analysed size for the participant-average estimand and unweighted for the
cluster-average estimand \citep{hayes2017, kahan2023}. This two-stage analysis
gives up some efficiency but its type I error rate is correct by construction
with any number of clusters.

\emph{Exact randomisation test.} Because randomisation acts on practices, the
set of allocations that could have occurred is finite and small: with 5 of 14
practices allocated to the intervention there are exactly
$\binom{14}{5} = 2002$ of them, and with a 7 versus 7 split, 3432. We will
therefore enumerate the full set, recompute the cluster-level test statistic
under each allocation, and take the proportion of allocations giving a
statistic at least as extreme as the observed one as an exact $p$ value. This
test makes no distributional assumption and is valid however few clusters there
are \citep{gail1996, leyrat2018}. Its resolution is limited by the design: the
smallest attainable two-sided $p$ value with a 5 versus 9 split is
$2/2002 \approx 0.001$. If randomisation was covariate-constrained, the
enumeration will be restricted to the allocations that satisfy the same
constraints, since only those could have been selected \citep{li2016}.

If the three analyses agree, the model-based estimate is reported as the
primary result. If they disagree, the disagreement itself will be reported and
the randomisation test treated as the arbiter of the significance claim, since
it is the only one of the three whose validity does not depend on the number of
clusters.

\subsubsection{Intra-cluster correlation}

We will report the estimated intra-cluster correlation for the primary outcome,
$\hat\rho = \hat\tau^2 / (\hat\tau^2 + \hat\sigma^2)$, with a 95\% confidence
interval obtained by profile likelihood, and, if the interval is very wide, we
will say so rather than presenting the point estimate alone. Reporting the
observed intra-cluster correlation is a CONSORT requirement for cluster trials
and is what makes the next trial in this field easier to design
\citep{campbell2012, adams2004}.

\subsection{Supportive analysis over both time points}

To keep the results comparable with the Australian trial \citep{bagg2021}, we
will additionally fit a mixed model for repeated measures using all
measurement occasions, with fixed effects for arm, time, and their interaction,
a random intercept for practice and a random intercept for participant nested
within practice. The treatment effect at 18 weeks is the corresponding
arm-by-time contrast. This analysis is supportive; \eqref{eq:primary} remains
the primary one. Mixed models for repeated measures behave acceptably in
cluster randomised trials with continuous outcomes missing at random
\citep{bell2020}.

\subsection{Secondary outcomes}

Secondary outcomes will be analysed with models of the same structure as
\eqref{eq:primary}, with the link function and distribution appropriate to the
outcome, always including a random intercept for practice and always using the
same small-sample correction. Effects will be reported as point estimates with
95\% confidence intervals (Table~\ref{tab:sec}). No formal claim of efficacy
will be based on a secondary outcome.

\subsection{Sensitivity analyses}

The following are prespecified and will be reported whether or not they change
the conclusion.

\begin{enumerate}[leftmargin=1.4em, itemsep=2pt]
\item \textbf{Leave-one-practice-out.} The primary model refitted 14 times,
each time omitting one practice. With few clusters a single atypical practice
can carry the result, and the spread of the 14 estimates shows how much it
does.
\item \textbf{Therapist effects.} A third level for therapist within practice.
This is the more faithful description of the data hierarchy, but variance
components at three levels are poorly identified with 14 practices, so it is a
sensitivity analysis rather than the primary model.
\item \textbf{Unadjusted analysis.} Model \eqref{eq:primary} without the
baseline covariate, to show what the baseline adjustment contributes.
\item \textbf{Cluster-mean baseline.} Model \eqref{eq:primary} with the
practice mean baseline score added as a further fixed effect, which separates
the within-practice from the between-practice association with baseline
disability.
\item \textbf{Wave.} If practices entered in distinct recruitment waves, the
primary model with wave as a fixed effect.
\end{enumerate}

\subsection{Missing data}

The primary analysis uses all participants with a recorded 18-week outcome and
is valid under a missing-at-random assumption conditional on the baseline score
and the covariates in the model. We will report the number and proportion of
missing outcomes by arm and by practice, and compare the baseline
characteristics of participants with and without an 18-week outcome. As a
sensitivity analysis we will use multiple imputation by chained equations with
the practice included in the imputation model, so that the imputation respects
the clustering \citep{bell2020}. If an entire practice is lost, that will be
reported separately and not imputed.

\subsection{Non-adherence and contamination}

Adherence will be summarised as the number of intervention sessions attended
and as the proportion of participants attending at least \open{threshold} of
them. If non-adherence is substantial, we will estimate the complier-average
causal effect using allocation as an instrument. With 14 clusters this analysis
is imprecise, and it is prespecified as exploratory.

Contamination --- control practices adopting elements of the intervention --- is
a real possibility in a trial where practices know each other. \open{describe
how contamination is measured, e.g. a therapist questionnaire at the end of the
trial}. Any contamination detected will be described and its likely direction
of effect discussed; it biases the estimated effect towards the null
\citep{hemming2021}.

\subsection{Interim analysis}
\label{sec:interim}

An interim analysis is planned and was specified in Bayesian terms for the
ethics submission. It is not described in this paper.
\open{decide: brief pointer here plus a companion paper, or a full section
added before submission}

% ============================================================
\section{Reporting}

Results will be reported following CONSORT 2025 with the extension for cluster
randomised trials \citep{hopewell2025, campbell2012}, including a flow diagram
that tracks both practices and participants (Figure~\ref{fig:flow}). Two items
of the cluster extension have shaped this plan directly: items 7a and 17a
require the intra-cluster correlation to be reported, with an indication of its
uncertainty, at the design stage and for each primary outcome; item 13a requires
the flow diagram to give the number of clusters as well as the number of
individuals at every stage \citep{campbell2012}.

The guideline for the content of statistical analysis plans that this document
follows \citep{gamble2017} contains no items specific to clustering; an
extension for cluster randomised trials is under development but not yet
published \citep{hemming2025sap}. We have therefore added the cluster-specific
items ourselves: the estimand and its weighting, the small-sample correction,
the intra-cluster correlation, and the two-level description of the sample.

% ============================================================
\section*{Declarations}

\textbf{Funding.} \open{funder and grant number}

\textbf{Conflicts of interest.} \open{to be completed by all authors}

\textbf{Ethics.} \open{committee, approval number, date}

\textbf{Data and code availability.} The code that reproduces the sample size
and power calculations reported here accompanies this paper. Analysis code for
the trial will be published together with the trial results.
\open{repository URL}

\textbf{Author contributions.} \open{CRediT statement}

% ============================================================
\clearpage
\section*{Tables and figures}

\begin{table}[ht]
\centering\small
\caption{Sample size per arm required to detect a 2-point difference in RMDQ at
18 weeks with 80\% power and two-sided $\alpha = 0.05$, by intra-cluster
correlation. Based on a pooled SD of 5.2 points, adjustment for the baseline
score (assumed correlation 0.6), 14 analysed participants per practice and a
coefficient of variation of practice size of 0.65.}
\label{tab:ss}
\begin{tabular}{lcccc}
\toprule
& \multicolumn{4}{c}{Intra-cluster correlation}\\
\cmidrule(l){2-5}
& 0 & 0.01 & 0.03 & 0.05\\
\midrule
Design effect (eq.~\ref{eq:de}) & 1.00 & 1.19 & 1.56 & 1.93\\
Analysed participants per arm   & 69   & 82   & 107  & 133\\
Recruited per arm (10\% attrition) & 76 & 91 & 119 & 148\\
\bottomrule
\end{tabular}
\end{table}

\begin{table}[ht]
\centering\small
\caption{Power to detect a 2-point difference in RMDQ at 18 weeks, two-sided
$\alpha = 0.05$, with a $t$ reference distribution on $k_1+k_2-2$ degrees of
freedom. Upper panel: the existing 14 practices with 200 analysed participants
in total. Lower panel: balanced designs with 14 analysed participants per
practice, so that the total sample grows with the number of practices.}
\label{tab:power}
\begin{tabular}{llccc}
\toprule
& & \multicolumn{3}{c}{Intra-cluster correlation}\\
\cmidrule(l){3-5}
Practices (int./ctrl.) & Analysed $N$ & 0.01 & 0.03 & 0.05\\
\midrule
\multicolumn{5}{l}{\emph{14 practices, $N$ fixed at 200}}\\
5 / 9  & 200 & 0.77 & 0.65 & 0.56\\
6 / 8  & 200 & 0.80 & 0.68 & 0.59\\
7 / 7  & 200 & 0.80 & 0.69 & 0.60\\
\midrule
\multicolumn{5}{l}{\emph{Balanced, 14 participants per practice}}\\
5 / 5   & 140 & 0.62 & 0.51 & 0.43\\
6 / 6   & 168 & 0.73 & 0.61 & 0.52\\
7 / 7   & 196 & 0.80 & 0.69 & 0.60\\
8 / 8   & 224 & 0.86 & 0.76 & 0.66\\
9 / 9   & 252 & 0.90 & 0.81 & 0.72\\
10 / 10 & 280 & 0.93 & 0.86 & 0.77\\
\bottomrule
\end{tabular}
\end{table}

\begin{table}[ht]
\centering\small
\caption{Baseline characteristics (shell). Panel A describes the clusters,
panel B the participants.}
\label{tab:base}
\begin{tabularx}{\textwidth}{Xccc}
\toprule
Characteristic & Intervention & Control & All\\
\midrule
\multicolumn{4}{l}{\emph{Panel A --- practices}}\\
Number of practices & xx & xx & xx\\
Therapists per practice, median (IQR) & xx (xx to xx) & xx (xx to xx) & xx (xx to xx)\\
Participants per practice, median (IQR) & xx (xx to xx) & xx (xx to xx) & xx (xx to xx)\\
Urban setting, n (\%) & xx (xx) & xx (xx) & xx (xx)\\
\midrule
\multicolumn{4}{l}{\emph{Panel B --- participants}}\\
Number of participants & xx & xx & xx\\
Age, mean (SD) & xx (xx) & xx (xx) & xx (xx)\\
Female, n (\%) & xx (xx) & xx (xx) & xx (xx)\\
Duration of current episode & xx (xx) & xx (xx) & xx (xx)\\
RMDQ, mean (SD) & xx (xx) & xx (xx) & xx (xx)\\
Pain intensity, mean (SD) & xx (xx) & xx (xx) & xx (xx)\\
\open{complete} & & &\\
\bottomrule
\end{tabularx}
\end{table}

\begin{table}[ht]
\centering\small
\caption{Analysis of the primary outcome (shell). All three analyses estimate
the participant-average difference in RMDQ at 18 weeks.}
\label{tab:prim}
\begin{tabularx}{\textwidth}{Xcccc}
\toprule
Analysis & Intervention & Control & Difference (95\% CI) & $p$\\
& mean (SD) & mean (SD) & &\\
\midrule
Mixed model, Kenward--Roger (primary) & xx (xx) & xx (xx) & xx (xx to xx) & .xx\\
Cluster-level, size-weighted & --- & --- & xx (xx to xx) & .xx\\
Exact randomisation test & --- & --- & xx (xx to xx) & .xx\\
Cluster-average effect (unweighted) & --- & --- & xx (xx to xx) & .xx\\
\midrule
Intra-cluster correlation (95\% CI) & \multicolumn{4}{l}{xx (xx to xx)}\\
\bottomrule
\end{tabularx}
\end{table}

\begin{table}[ht]
\centering\small
\caption{Analysis of secondary outcomes (shell).}
\label{tab:sec}
\begin{tabularx}{\textwidth}{Xcccc}
\toprule
Outcome & Intervention & Control & Effect (95\% CI) & $p$\\
\midrule
\open{list} & xx (xx) & xx (xx) & xx (xx to xx) & .xx\\
\bottomrule
\end{tabularx}
\end{table}

\begin{figure}[ht]
\centering
\fbox{\parbox{0.9\textwidth}{\centering\vspace{3.2cm}
CONSORT flow diagram for cluster randomised trials (shell): practices
screened, randomised and analysed in the upper tier; participants screened,
recruited, followed up and analysed in the lower tier, within practices.
\open{draw}\vspace{3.2cm}}}
\caption{Flow of practices and participants through the trial.}
\label{fig:flow}
\end{figure}

% ============================================================
\clearpage
\bibliographystyle{unsrtnat}
\bibliography{references}

\end{document}
